\section{Learning executable monodromy classes}
\label{sec:monodromy-principle}

The mechanism studied here separates two operations that are easily conflated:
changing the implementation of an existing computation, and acquiring a
different executable computation. A deformation that is certified to preserve
a based-loop homotopy class can support the former, but cannot by itself
accomplish the latter. We therefore develop a learning interface with
class-changing proposals for unadmitted operations and class-preserving
transactions for admitted operations. The current evidence establishes this
interface in a small supplied covering; it does not yet establish an independent
topological advantage over equally informed conventional optimization.

\subsection{Covering space, marked states, and executable composition}

For an integer $D\geq 2$, consider the supplied polynomial family
\begin{equation}
 F_D(z,u)=z^D-Dz-u,
 \qquad
 \mathcal C_D=\{-(D-1)\zeta:\zeta^{D-1}=1\},
 \qquad U_D=\mathbb C\setminus\mathcal C_D.
 \label{eq:cover-and-critical-set}
\end{equation}
Projection onto $u$ defines an unbranched $D$-sheet covering over $U_D$.
We fix the control basepoint $u_*=0$ and mark its $D$ distinct roots once.
Lifting a based loop $\gamma:[0,1]\to U_D$ transports each marked root to
another root in the same fiber. Its endpoint action is a permutation
$\Pi_\gamma=\rho_D([\gamma])$. For chronological concatenation, defined so
that $\gamma_1\star\gamma_2$ executes $\gamma_1$ first,
\begin{equation}
 \Pi_{\gamma_1\star\gamma_2}
   =\Pi_{\gamma_2}\circ\Pi_{\gamma_1}.
 \label{eq:chronological-composition}
\end{equation}
An internal state is a sheet index, not a freshly initialized root at every
call. Equation~\eqref{eq:chronological-composition} consequently specifies an
executable, generally noncommutative composition rule.

The executed class-acquisition diagnostic uses $D=3$. The roots at $u_*=0$
are labelled $0,1,2$ in the order $(-\sqrt{3},0,\sqrt{3})$, and the critical
controls are $-2$ and $2$. Its one-parameter family is the polygonal loop
\begin{equation}
 \gamma_a:0\longrightarrow a+i\longrightarrow a-i\longrightarrow0,
 \qquad
 \Pi_{\gamma_a}=
 \begin{cases}
 B=(12),&a<-2,\\
 I,&-2<a<2,\\
 A=(01),&a>2.
 \end{cases}
 \label{eq:cubic-triangle-chambers}
\end{equation}
The singular cases $a=\pm2$ are not executable. For example,
$\Pi_B(\Pi_A(0))=2$, whereas $\Pi_A(\Pi_B(0))=1$; the order of calls matters.
The permutation labels in \eqref{eq:cubic-triangle-chambers} describe the
independently checked outcome, rather than an operation table provided to the
learner. The numerical executor actually continues all three roots with a
safeguarded predictor/Newton corrector. Its full-fiber endpoint observations,
known covering, critical controls, and numerical costs are available to every
comparison arm. Thus the protocol combines terminal task labels with
full-fiber execution observations; it is not a single-sheet-only observation
protocol. The executed source and independent reconstruction are preserved in
\path{.deep-work/2026-10-31T234403-0400-v2/probes/topology-architecture/versions/code-parameter-loop-boundary-learning-v1/loop_boundary_learning.py}
and
\path{.deep-work/2026-10-31T234403-0400-v2/supervision/reviews/review-038-loop-boundary-learning.md}.

\subsection{What safe implementation deformation does and does not learn}

An earlier construction first identifies a permutation, compiles it into a
loop of distinct root configurations, and then changes its geometry using a
time-dependent ambient shear. Writing a point as $x+iy$, this shear is
\begin{equation}
 \Phi_{\omega,t}(x+iy)
 =x+i\left[y+\sin^2(\pi t/T)h_\omega(x)\right],
 \qquad h_\omega(x)=\sum_\ell\omega_\ell\tanh(v_\ell x+b_\ell),
 \label{eq:safe-ambient-shear}
\end{equation}
where $v_\ell,b_\ell$ are fixed. Every finite $\omega$ gives a global
diffeomorphism, with an explicit inverse obtained by subtracting the same
vertical displacement. At both execution endpoints it is the identity.
Interpolating $\omega$ therefore gives a collision-free homotopy of
\emph{configuration-space} loops with a fixed base configuration, preserving
their braid and endpoint permutation. This is a useful within-class learning
freedom, but it is not a mechanism for discovering a new permutation. It is
also a different construction from the control-plane covering in
\eqref{eq:cover-and-critical-set}; the two spaces must not be identified
without an additional map between them.

The new class-learning interface instead trains the raw vertices of a loop in
the fixed control plane. A controller $a_\theta(e)$ maps a public operation
feature $e$ to coordinates in a supplied affine polygon family. For an
unadmitted operation, a finite update $\theta\mapsto\theta^+$ may place the
two regular endpoint snapshots in different chambers. The parameter
interpolation between them may pass through a singular loop, but that
interpolation is not executed. Only the separately checked regular snapshots
are run. There is consequently no claim that a nonsingular based homotopy can
change monodromy. The covering, basepoint, critical set and sheet labels remain
fixed. Moreover, a change of homotopy class in the kernel of $\rho_D$ would
not count as a new executable operation.

\subsection{Terminal boundary feedback and neural credit}

An exact endpoint-semantic loss is locally constant in regular loop
coordinates. Ordinary within-chamber differentiation alone therefore provides
no direct derivative for acquiring a different endpoint action. Let $V$ be
raw loop coordinates drawn from a smooth location kernel centered at
$a_\theta$, and let $\ell(V;\theta)$ denote the terminal loss obtained by
executing that loop and the current neural bindings. Define
\begin{equation}
 J_\sigma(\theta)
 =\int \ell(v;\theta)p_\sigma(v-a_\theta)\,dv.
 \label{eq:coupled-smoothed-objective}
\end{equation}
For the current fixed-cover, endpoint-semantic objective, the coupled
derivative has two contributions:
\begin{equation}
 \nabla_\theta J_\sigma
 =\mathbb E\!\left[\left.\partial_\theta\ell(V;\theta)
                    \right|_{V\ \mathrm{fixed}}\right]
 +(D_\theta a_\theta)^\top g_{\mathrm{boundary}}.
 \label{eq:coupled-boundary-derivative}
\end{equation}
The first term trains the input encoding, readout, calls and their shared
backbone while holding the actually observed endpoint maps fixed. For a
coordinate slice with crossing events $b_r$ and terminal-loss jumps
$\Delta\ell_r=\ell(b_r^+;\theta)-\ell(b_r^-;\theta)$, the second term uses
\begin{equation}
 (g_{\mathrm{boundary}})_j
 =\mathbb E_{V_{-j}}\!\left[
       \sum_r\Delta\ell_r\,p_\sigma(b_r-a_{\theta,j})\right].
 \label{eq:boundary-jump-estimator}
\end{equation}
Safe representatives of the neighboring chambers are actually executed to
obtain these terminal losses. The implementation uses a normalized
raised-cosine kernel and charges event detection, representatives, terminal
callbacks and continuation. A coordinate moving one free vertex affects its
two adjacent segments; tied coordinates moving several vertices are charged
for every affected segment instead of silently inheriting the smaller bound.

Equations~\eqref{eq:coupled-boundary-derivative}--\eqref{eq:boundary-jump-estimator}
are applications of known differentiation of piecewise-smooth objectives,
not a claimed new calculus theorem. A path-work objective with smooth
within-chamber dependence on $v$ would require its corresponding continuous
derivative as well. The implemented boundary term changes how probability
mass is assigned to executable classes; it does not replace successful
discrete execution by a mixture of several operations.

The ideal smoothing distribution assigns mass arbitrarily close to a
discriminant. The numerical executor additionally requires positive finite
clearance. These are distinct domains. The initial score-function arm was
censored by a probe entering the exclusion band, rather than assigned a
fabricated failure loss. The later hybrid interface evaluates an ideal
sample's class at a paid safe representative when available, while still
requiring the deterministic candidate loop itself to be executable. This
scalar-family convention is exposed to the conventional controls and is not
assumed to solve the general high-dimensional feasibility problem.

\subsection{A sufficient protected transaction}

For an admitted polygonal loop $\gamma$, let
$\delta=\min_{t,c\in\mathcal C_D}|\gamma(t)-c|>0$. If its updated vertices
satisfy
\begin{equation}
 \max_k|v_k(\theta^+,e)-v_k(\theta,e)|<\delta,
 \label{eq:protected-vertex-transaction}
\end{equation}
with the same segmentation and fixed endpoints, linear interpolation gives a
based homotopy avoiding $\mathcal C_D$. Indeed, the displacement of every
point on every interpolated segment is bounded by the maximum vertex
displacement; its distance from the critical set remains positive. Hence its
endpoint permutation is unchanged. Preserving each admitted generator then
preserves every word of those generators under fixed input/output bindings.

This sufficient topological certificate is accompanied by separate checks of
the executor's coordinate box and numerical clearance. In the diagnostic the
minimum allowed control clearance is $.015$. Protection of a permutation
alone does not protect a changing encoder, router or decoder; the additional
binding and soft-execution contracts are specified in
Section~\ref{sec:prometheus-block}. None of these conditions implies hardware
noise protection or a quantum energy advantage from a CPU simulation.

\subsection{Executed class acquisition and the unresolved advantage}

The initial scalar diagnostic uses two real shared weights,
\begin{equation}
 a_{\mathrm{old}}=w_0,\qquad
 a_{\mathrm{new}}=w_0+w_1,\qquad (w_0,w_1)=(-3,4).
\end{equation}
The old operation $B$ is supplied as the starting mechanism and verified by
continuation; it is not credited as learned in an earlier stage. Four
terminal training cases require $A$, which is outside the admitted subgroup
$\{I,B\}$. Two accepted neural updates produce
\begin{center}
\begin{tabular}{c|rr|rr}
snapshot & $a_{\mathrm{old}}$ & $a_{\mathrm{new}}$
         & training error & old error\\\hline
initial & $-3$ & $1$ & $.75$ & $0$\\
first & $-2.625$ & $1.75$ & $.75$ & $0$\\
second & $-2.26427259$ & $2.47145482$ & $0$ & $0$
\end{tabular}
\end{center}
The old loop actually changes geometry while retaining $B$; the new loop
crosses its class boundary and acquires $A$. Four frozen held-out composition
cases and an actual checkpoint reload also have zero error. These are finite
length-at-most-three tests, not a scaling or long-term continual-learning
result. The raw execution is
\path{.deep-work/2026-10-31T234403-0400-v2/probes/topology-architecture/runs/parameter-loop-boundary-learning-v1/learn/result.json}.

The equally informed ordinary event optimizer produces the same weights,
parameter hashes and counted trajectory. Discrete chamber search succeeds
with one update, $35$ program executions and $37$ terminal callbacks, versus
$70$ executions and $72$ callbacks for the event arm. Those totals include
their disclosed calibration and audit work. A reexecution guard spends more
continuation work than the analytic certificate on this instance, but the
ordinary optimizer is also allowed that certificate. Thus the result closes
the earlier map-first learning gap, while the principal question remains open:
does monodromy organization improve acquisition, retention, verification or
parameter growth beyond the strongest ordinary constraint and search methods
under matched information and complete cost?
